\section{Introduction}
The Agile Manifesto revolutionized software engineering by prioritizing
working software, individual interactions, and responding to change over
comprehensive documentation and strict upfront planning. For over two decades,
this empirical paradigm was the gold standard for navigating requirements
uncertainty, and for good reason: in a world where local implementation effort
and changing requirements made heavyweight plans costly, minimizing unnecessary
upfront documentation helped reduce waste.

That world is changing. Generative AI assistants now produce substantial
amounts of code in seconds, shifting scarce software-engineering effort toward
architectural alignment, integration, and cognitive coordination. Under this
new cost equation, the Agile bet---that planning can remain deliberately light
as understanding develops---deserves to be tested rather than assumed.

If deferring design in an AI-assisted environment merely relocates effort, late
project activity may reflect several mechanisms: compensatory recovery, delayed
integration without assessed improvement, distributed progress, or
technical-complexity pressure. This paper examines how such mechanisms become
visible through repository metadata, multi-checkpoint student surveys,
independent evaluator scorecards, and qualitative meeting transcripts across
two cohorts of software-engineering projects. The resulting evidence connects
what students report, what repositories reveal, what evaluators score, and what
coordination language appears in the project record.

We structure the investigation around three questions:

\begin{itemize}
	\item \textbf{RQ1 (Descriptive):} How do student developers utilize Generative AI, and how does this adoption shape their longitudinal perceptions of software engineering roles and project expectations?
	\item \textbf{RQ2 (Descriptive):} How does the temporal density of repository activity contrast with the qualitative typification of human coordination friction across the project lifecycle?
	\item \textbf{RQ3 (Primary):} How does the observable scope and content evidence of upfront architectural artifacts at baseline $T_1$ associate with late-stage clean-path rework and independent, blind-scored project outcomes?
\end{itemize}

The study makes three contributions. First, it provides a longitudinal,
multi-method account of AI-assisted student development that keeps perceived
AI impact, self-reported adoption, repository activity, and qualitative
coordination evidence analytically distinct. Second, it makes the temporal
contrast between repository activity and integration pressure observable and
qualifies that contrast using non-overlapping phase bins, timing/outcome
contrasts, operational regularity, and coverage-aware transcript evidence while
preserving the different grains and coverage of the underlying evidence.
Third, it examines upfront architectural-artifact evidence, late clean-path
rework, and independent outcomes without collapsing structural observations
into an unsupported quality score. The resulting picture is deliberately more
useful than a simple success/failure verdict: rather than eliminating threats
such as ordinary deadline pressure, process heterogeneity, small-sample
sensitivity, and technical complexity, it makes them explicit and inspectable.
It shows where implementation speed meets the harder problem of shared system
intent.

The findings motivate a practical proposition. Generative AI does not remove
the need for Agile practice; it increases the importance of making architectural
intent, interfaces, dependencies, and acceptance conditions explicit while
remaining open to revision. These observations motivate Spec-Driven
Development (SDD) as a theoretically grounded, testable hypothesis: carrying a
specification through technical planning, task decomposition, implementation,
and convergence checks may help teams address the coordination bottlenecks
exposed by the combined timing, planning, rework, complexity, and influence
diagnostics \cite{github2025spec,thoughtworks202Xspec}. This observational
study neither implemented SDD nor compared it with another development
approach; therefore, it does not establish SDD's effectiveness. Instead, the
hypothesis invites future intervention research into how teams can preserve
adaptability while ensuring that people and AI systems act on the same
inspectable account of what the system is meant to become.
